% GENERATED FILE. Edit canonical lessons, then run npm run generate:books.
% canonical-source-hash: fnv1a64:90df147afd279212
% canonical-lessons: KA-C459-parvata, KA-C459-sikhara, KA-C459-muggarisu, KA-C459-dere, KA-C459-bidara

\chapter{Mountains and Camps}
\label{ch:ka-mountains-and-camps}
\hlchaptermodality{459}

\begin{chapteropening}
\textbf{By the end of this chapter:} \emph{I can say a mountain, a peak, to stumble, a tent and an encampment.}

Complete the last lesson of chapter 459: I can say a mountain, a peak, to stumble, a tent and an encampment.
\end{chapteropening}

\section[\texorpdfstring{\kn{ಪರ್ವತ} (parvata)}{ಪರ್ವತ (parvata)}]{\kn{ಪರ್ವತ} (parvata) — a mountain}

\label{lesson:KA-C459-parvata}



\begin{quote}
Before the new one: say the Kannada for to get back, then the Kannada for to stay overnight.
\end{quote}

\subsection*{You'll want to know: \kn{ಪರ್ವತ}}
\textbf{\kn{ಪರ್ವತ}} — \emph{parvata} — \textquotedblleft{}a mountain\textquotedblright{}.

\begin{quote}
\kn{ಆ} \kn{ಪರ್ವತದ} \kn{ಮೇಲೆ} \kn{ಯಾವಾಗಲೂ} \kn{ಹಿಮ} \kn{ಇರುತ್ತದೆ}.
\end{quote}

\begin{quote}
\emph{ā parvatada mēle yāvāgalū hima iruttade.}
\end{quote}

\begin{quote}
There is always snow on that mountain.
\end{quote}

\kn{ಬೆಟ್ಟ} is a hill; \kn{ಪರ್ವತ} is a high mountain.

\begin{cousinweb}
\textbf{\kn{ಪರ್ವತ}} (\emph{parvata}) is borrowed from Sanskrit \emph{parvata}, \textquotedblleft{}a rock, a mountain\textquotedblright{} (CDIAL 7945).
\end{cousinweb}

\begin{grammarlens}[title={What you've built}]
One more word for getting around.
\end{grammarlens}

\subsection*{Your turn}
\begin{itemize}
\raggedright
  \item \emph{Say it:} \emph{parvata}
  \item \emph{Say it:} \emph{parvata}, once more
  \item \emph{Say it:} \emph{ā parvatada mēle yāvāgalū hima iruttade.} — There is always snow on that mountain.
  \item \emph{Recall:} say the Kannada for to get back, then the Kannada for to stay overnight, then say \emph{parvata} again
\end{itemize}

\begin{tcolorbox}[breakable,colback=teal!4,colframe=teal!35!black,title={Before you move on}]
What does \kn{ಪರ್ವತ} mean? (\textquotedblleft{}A mountain\textquotedblright{}.) Say it once more.
\end{tcolorbox}
\section[\texorpdfstring{\kn{ಶಿಖರ} (śikhara)}{ಶಿಖರ (śikhara)}]{\kn{ಶಿಖರ} (śikhara) — a peak; a summit}

\label{lesson:KA-C459-sikhara}



\begin{quote}
Before the new one: say the Kannada for to stay overnight, then the Kannada for a mountain.
\end{quote}

\subsection*{You'll want to know: \kn{ಶಿಖರ}}
\textbf{\kn{ಶಿಖರ}} — \emph{śikhara} — \textquotedblleft{}a peak; a summit\textquotedblright{}.

\begin{quote}
\kn{ನಾವು} \kn{ಬೆಳಗ್ಗೆ} \kn{ಶಿಖರ} \kn{ತಲುಪಿದೆವು}.
\end{quote}

\begin{quote}
\emph{nāvu beḷagge śikhara talupidevu.}
\end{quote}

\begin{quote}
We reached the summit in the morning.
\end{quote}

The top of a mountain, and the top of a temple tower.

\begin{cousinweb}
\textbf{\kn{ಶಿಖರ}} (\emph{śikhara}) is borrowed from Sanskrit \emph{śikhara}, \textquotedblleft{}a point, a peak\textquotedblright{} (CDIAL 12435).
\end{cousinweb}

\begin{grammarlens}[title={What you've built}]
One more word for getting around.
\end{grammarlens}

\subsection*{Your turn}
\begin{itemize}
\raggedright
  \item \emph{Say it:} \emph{śikhara}
  \item \emph{Say it:} \emph{śikhara}, once more
  \item \emph{Say it:} \emph{nāvu beḷagge śikhara talupidevu.} — We reached the summit in the morning.
  \item \emph{Recall:} say the Kannada for to stay overnight, then the Kannada for a mountain, then say \emph{śikhara} again
\end{itemize}

\begin{tcolorbox}[breakable,colback=teal!4,colframe=teal!35!black,title={Before you move on}]
What does \kn{ಶಿಖರ} mean? (\textquotedblleft{}A peak; a summit\textquotedblright{}.) Say it once more.
\end{tcolorbox}
\section[\texorpdfstring{\kn{ಮುಗ್ಗರಿಸು} (muggarisu)}{ಮುಗ್ಗರಿಸು (muggarisu)}]{\kn{ಮುಗ್ಗರಿಸು} (muggarisu) — to stumble; to lurch}

\label{lesson:KA-C459-muggarisu}



\begin{quote}
Before the new one: say the Kannada for a mountain, then the Kannada for a peak.
\end{quote}

\subsection*{You'll want to know: \kn{ಮುಗ್ಗರಿಸು}}
\textbf{\kn{ಮುಗ್ಗರಿಸು}} — \emph{muggarisu} — \textquotedblleft{}to stumble; to lurch\textquotedblright{}.

\begin{quote}
\kn{ಕಲ್ಲಿನ} \kn{ದಾರಿಯಲ್ಲಿ} \kn{ಅಜ್ಜ} \kn{ಮುಗ್ಗರಿಸಿದರು}.
\end{quote}

\begin{quote}
\emph{kallina dāriyalli ajja muggarisidaru.}
\end{quote}

\begin{quote}
Grandfather stumbled on the stony path.
\end{quote}

Past (I …): \kn{ಮುಗ್ಗರಿಸಿದೆ}, \emph{muggariside}.

A heavier stumble than \kn{ಎಡವು}, to trip.

\begin{cousinweb}
\textbf{\kn{ಮುಗ್ಗರಿಸು}} (\emph{muggarisu}) is a Kannada verb whose origin is not settled; DEDR shows no clear source for it.
\end{cousinweb}

\begin{grammarlens}[title={What you've built}]
One more word for getting around.
\end{grammarlens}

\subsection*{Your turn}
\begin{itemize}
\raggedright
  \item \emph{Say it:} \emph{muggarisu}
  \item \emph{Say it:} \emph{muggarisu}, once more
  \item \emph{Say it:} \emph{kallina dāriyalli ajja muggarisidaru.} — Grandfather stumbled on the stony path.
  \item \emph{Recall:} say the Kannada for a mountain, then the Kannada for a peak, then say \emph{muggarisu} again
\end{itemize}

\begin{tcolorbox}[breakable,colback=teal!4,colframe=teal!35!black,title={Before you move on}]
What does \kn{ಮುಗ್ಗರಿಸು} mean? (\textquotedblleft{}To stumble; to lurch\textquotedblright{}.) Say it once more.
\end{tcolorbox}
\section[\texorpdfstring{\kn{ಡೇರೆ} (ḍēre)}{ಡೇರೆ (ḍēre)}]{\kn{ಡೇರೆ} (ḍēre) — a tent}

\label{lesson:KA-C459-dere}



\begin{quote}
Before the new one: say the Kannada for a peak, then the Kannada for to stumble.
\end{quote}

\subsection*{You'll want to know: \kn{ಡೇರೆ}}
\textbf{\kn{ಡೇರೆ}} — \emph{ḍēre} — \textquotedblleft{}a tent\textquotedblright{}.

\begin{quote}
\kn{ರಾತ್ರಿ} \kn{ನಾವು} \kn{ಶಿಖರದ} \kn{ಕೆಳಗೆ} \kn{ಡೇರೆ} \kn{ಹಾಕಿದೆವು}.
\end{quote}

\begin{quote}
\emph{rātri nāvu śikharada keḷage ḍēre hākidevu.}
\end{quote}

\begin{quote}
At night we pitched a tent below the summit.
\end{quote}

\kn{ಡೇರೆ} \kn{ಹಾಕು}, to pitch a tent.

\begin{cousinweb}
\textbf{\kn{ಡೇರೆ}} (\emph{ḍēre}) is usually traced to Persian \emph{dēra}, \textquotedblleft{}a tent, a camp\textquotedblright{}, which reached Kannada through Hindustani.
\end{cousinweb}

\begin{grammarlens}[title={What you've built}]
One more word for getting around.
\end{grammarlens}

\subsection*{Your turn}
\begin{itemize}
\raggedright
  \item \emph{Say it:} \emph{ḍēre}
  \item \emph{Say it:} \emph{ḍēre}, once more
  \item \emph{Say it:} \emph{rātri nāvu śikharada keḷage ḍēre hākidevu.} — At night we pitched a tent below the summit.
  \item \emph{Recall:} say the Kannada for a peak, then the Kannada for to stumble, then say \emph{ḍēre} again
\end{itemize}

\begin{tcolorbox}[breakable,colback=teal!4,colframe=teal!35!black,title={Before you move on}]
What does \kn{ಡೇರೆ} mean? (\textquotedblleft{}A tent\textquotedblright{}.) Say it once more.
\end{tcolorbox}
\section[\texorpdfstring{\kn{ಬಿಡಾರ} (biḍāra)}{ಬಿಡಾರ (biḍāra)}]{\kn{ಬಿಡಾರ} (biḍāra) — an encampment; a camp}

\label{lesson:KA-C459-bidara}



\begin{quote}
Before the new one: say the Kannada for to stumble, then the Kannada for a tent.
\end{quote}

\subsection*{You'll want to know: \kn{ಬಿಡಾರ}}
\textbf{\kn{ಬಿಡಾರ}} — \emph{biḍāra} — \textquotedblleft{}an encampment; a camp\textquotedblright{}.

\begin{quote}
\kn{ಕೂಲಿಗಳ} \kn{ಬಿಡಾರ} \kn{ನದಿಯ} \kn{ದಡದಲ್ಲಿ} \kn{ಇತ್ತು}.
\end{quote}

\begin{quote}
\emph{kūligaḷa biḍāra nadiya daḍadalli ittu.}
\end{quote}

\begin{quote}
The labourers’ camp was on the riverbank.
\end{quote}

A place of tents or huts put up for a while.

\begin{cousinweb}
\textbf{\kn{ಬಿಡಾರ}} (\emph{biḍāra}) is native Dravidian: DEDR 5393 lists Kannada \emph{biḍāra}, \textquotedblleft{}a halting place\textquotedblright{}, in the family of \kn{ಬಿಡು}, \textquotedblleft{}to let go; to stop, to halt\textquotedblright{}.
\end{cousinweb}

\begin{grammarlens}[title={What you've built}]
One more word for getting around.
\end{grammarlens}

\subsection*{Your turn}
\begin{itemize}
\raggedright
  \item \emph{Say it:} \emph{biḍāra}
  \item \emph{Say it:} \emph{biḍāra}, once more
  \item \emph{Say it:} \emph{kūligaḷa biḍāra nadiya daḍadalli ittu.} — The labourers’ camp was on the riverbank.
  \item \emph{Recall:} say the Kannada for to stumble, then the Kannada for a tent, then say \emph{biḍāra} again
\end{itemize}

\begin{tcolorbox}[breakable,colback=teal!4,colframe=teal!35!black,title={Before you move on}]
What does \kn{ಬಿಡಾರ} mean? (\textquotedblleft{}An encampment; a camp\textquotedblright{}.) Say it once more.
\end{tcolorbox}
